\documentclass[aspectratio=1610,onlymath]{beamer}
% \documentclass[aspectratio=1610,onlymath,handout]{beamer}

\input{macros-lecture}

\defineTitle{1}{Willkommen/Einleitung/Übersicht}{13. April 2026}

\begin{document}

\maketitle


\sectionSlide{Willkommen zur Vorlesung\\Theoretische Informatik und Logik\\{\large (ab Oktober 2026: ``Logik und Komplexität'')}}

% \begin{frame}\frametitle{}
%
% ~\hfill
% \includegraphics[height=6.5cm]{a1}
% \hfill~
%
% \end{frame}

\begin{frame}\frametitle{Raum, Zeit, URL}

\begin{itemize}
\item \emph{Vorlesungen:}\\
	montags DS3 (11:10–12:40) in HÜL/S386\\
	donnerstags DS4 (13:00–14:30) in HSZ/0002
\item \emph{Keine Vorlesungen:}\\
%	Feiertag: 1.5. (Tag der Arbeit)\\
	Feiertag: 14.5. (Himmelfahrt)\\
% 	Karfreitag: Fr 14.4.\\
% 	Dies academicus: Mi 17.5.\\
	Pfingstferien: Mo 25.5., Do 28.5.
\item \emph{Vorlesungswebseite}:\\[0.5ex]
	\url{\lectureurl}\\[0.5ex]
	(Folien, Übungsblätter, Termine, etc.)
\item \emph{Quellcode}:\\[0.5ex]
	\url{https://github.com/knowsys/TheoLog}\\[0.5ex]
	(Fragen, Bug-Reports, Pull-Requests sind willkommen)
\item \emph{Matrix (Chat)}:\\[0.5ex]
	\url{https://matrix.tu-dresden.de/\#/room/\#theolog:tu-dresden.de}\\[0.5ex]
	(Support bei allen Fragen zur Vorlesung)
\end{itemize}

\end{frame}


\begin{frame}\frametitle{Übungen}
\begin{itemize}
\item \emph{Anmeldung} zu den Übungen über OPAL \\(Link siehe Vorlesungswebseite; Beginn Einschreibung 13.4. um 15:00 Uhr)
\item \emph{Übungsblätter} jeweils freitags für übernächste Woche
% \\
% 	(erstes Übungsblatt am 13. Oktober 2016)
\item \emph{Beginn der Übungen:} 20. April 2026
\item \emph{Übungsablauf}, vereinfacht, idealisiert:\\
	Aufgaben werden zu Hause bearbeitet so gut es geht;\\
	in der Übung helfen Gruppenleiter bei Fragen und Problemen und zeigen Beispiellösungen
% 	\pause\\[1ex]
% \item \emph{Motivation:}\\
% 	Die Klausur am Semesterende wird zwei Aufgaben enthalten, die zuvor in der Übung
% 	besprochen wurden
\end{itemize}

\end{frame}


\begin{frame}\frametitle{Prüfung und Prüfungsvorbereitung}
\begin{itemize}
\item schritliche Prüfung (90min) am Ende des Sommersemesters
\item kann für alte und neue Studienordnungen belegt werden (INF-B-290, INF-D-330, INF-LE-EUI, INF-25-BA-LUK)
\item prüfungsrelevant:\\
	kompletter Stoff aus Vorlesung \alert{und} Übung;\\
	Wiedergeben (Definieren), Anwenden (Rechnen) und Erklären (Beweisen)
\item Modulnote ergibt sich je nach Studiengang
\item zur zusätzlichen Vorbereitung gibt es \alert{mehrere Konsultationen} und \alert{eine Probeklausur}
\end{itemize}

\end{frame}


\begin{frame}\frametitle{2. Wiederholungsprüfung TheoLog?}

\redalert{ACHTUNG!}
\begin{itemize}
\item Im Oktober 2026 findet ein automatischer Übertritt in die neuen Studienordnungen statt 
\item TheoLog bestanden? $\leadsto$ Modulnote wird für INF-25-BA-LUK übernommen
\item TheoLog nicht bestanden? $\leadsto$ Fehlversuche werden nicht übernommen
\item Wer Sommer eine 2. Wiederholungsprüfung antritt und nicht besteht, beendet sein Studium vor dem Übertritt. Der zweite Fehlversuch kann nicht gelöscht werden.
\item Für die 2. Wiederholung müssen Sie sich zum nächstmöglichen Termin (SoSe'26) anmelden, wenn (1) der schriftliche Bescheid zum ``wiederholten Nichtbestehen'' rechtzeitig zugestellt wurde und (2) Sie nicht in einem Urlaubssemester sind. Sie müssen die Prüfung dann auch ablegen (falls Sie nicht am Prüfungstag krank sind).
\end{itemize}
\redalert{$\leadsto$ Informieren Sie sich über die für Sie geltenden Fristen!} 
\medskip

{\tiny Analoge Bemerkungen gelten für viele andere Prüfungen im SoSe 2026.\\ Siehe \url{https://tu-dresden.de/ing/informatik/studium/examination-office/ueberleitung-zwangsuebertritt}

}

\end{frame}


% \begin{frame}\frametitle{TheoLog Bestehen}
% Tipps:
% \begin{itemize}
% \item \emph{Von Hand Mitschreiben}\\
% 	{\footnotesize Man merkt sich Stoff deutlich besser, wenn man ihn für sich selbst handschriftlich
% 	zusammenfasst.\footnote{\tiny P.\ Mueller \& D.\ Oppenheimer. The Pen Is Mightier Than the Keyboard: Advantages of Longhand Over Laptop Note Taking. Psychological Science, 06/2014, 25:6}}
% \item \emph{Selber Rechnen}\\
% 	{\footnotesize Die Prüfung besteht im Lösen von Rechenaufgaben. Theorie allein hilft da nicht weiter.}
% \item \emph{Schnell sein}\\
% 	{\footnotesize Prüfungszeit ist meistens knapp. Es reicht nicht, Aufgaben "`im Prinzip"' lösen zu können. Man muss sie schnell lösen.}
% \item \emph{Ehrlich zu sich selbst sein}\\
% 	{\footnotesize Man sollte selbst wissen, ob man genug gelernt hat oder nicht.\footnote{\tiny Vgl.\ aber auch Wikipedia \href{https://de.wikipedia.org/wiki/Dunning-Kruger-Effekt}{[[Dunning-Kruger-Effekt]]}}}
% \end{itemize}
%
% \end{frame}



% 
% \begin{frame}\frametitle{Kernfragen der theoretischen Informatik}
% 
% \begin{itemize}
% \item Was heißt "`berechnen"'?
% \item Welche Probleme kann man auf reellen Computern lösen?
% \item Was wäre wenn wir mächtigere Computer hätten?
% \item Was macht Rechenprobleme "`schwer"' oder "`einfach"'?
% \item Sind alternative Rechenmodelle denkbar?
% \item Welche (mathematischen/physikalischen/biologischen) Systeme können sonst noch rechnen?
% \end{itemize}
% 
% Diese finden sich wieder in zahlreichen Teilgebieten (Berechenbarkeit, Automaten, Komplexität, Quantencomputing, logisches Schließen, künstliche Intelligenz, \ldots)
% 
% \end{frame}

\begin{frame}\frametitle{Übersicht "`TheoLog"' (vorläufige Planung)}

\alert{Teil 1: Aussagenlogik}
\begin{itemize}
\item Syntax und Semantik
\item logisches Schließen (Resolution)
\item Horn-Logik als Vereinfachung
\item Praktische Anwendungen, SMT-Solver
\end{itemize}

\alert{Teil 2: Komplexität}
\begin{itemize}
\item Zeit und Raum
\item P und NP
\item PSpace
\end{itemize}

\alert{Teil 3: Prädikatenlogik}
\begin{itemize}
\item Syntax und Semantik
\item Modell-Checking (Anwendung: Datenbanken)
\item Logisches Schließen mit Resolution
\end{itemize}

\alert{Teil 4: Weitere Themen}
\begin{itemize}
\item Datalog
\item Modallogiken
\item Gödel (wenn die Zeit ausreicht)
% \item \ldots
\end{itemize}

\end{frame}

\begin{frame}\frametitle{Literatur: Lehrbücher}

Der Vorlesungsstoff gehört zu fast jeder Informatikausbildung. Es gibt viele
Lehrmaterialien und eine weitgehend einheitliche \ghost{Notation.}
\bigskip

\begin{itemize}
\item Uwe Schöning: \emph{Theoretische Informatik -- kurz gefasst.} Spektrum Akademischer Verlag\\
{\footnotesize\textcolor{devilscss}{klassischer deutschsprachiger Text}}
%
\item Michael Sipser: \emph{Introduction to the Theory of Computation.}  Cengage Learning\\
{\footnotesize\textcolor{devilscss}{Standardtext zu Sprachen und Berechnungskomplexität}}
%
\item Christopher Moore, Stephan Meterns: \emph{The Nature of Computation.} Oxford University Press\\
{\footnotesize\textcolor{devilscss}{moderner Text zu Komplexität und Berechnung, weniger formell}}
\end{itemize}

Zu speziellen Themen der Logik werden wir bei Gelegenheit noch Literaturangaben machen.

\end{frame}

\begin{frame}\frametitle{Literatur: Mehr Bücher}

Teile des Stoffs der Vorlesung werden auch in vielen, zum Teil recht kurzweiligen
Büchern weniger formal behandelt, z.B.:
\bigskip

\begin{itemize}
\item Apostolos Doxiadis, Christos Papadimitriou:\\ \emph{Logicomix: An Epic Search for Truth.} Bloomsbury\\
{\footnotesize\textcolor{devilscss}{Graphic Novel, inspiriert von Russels Leben und der Geschichte der \ghost{Logik}}}
%
\item Scott Aaronson: \emph{Quantum Computing Since Democritus.} Cambridge\\
{\footnotesize\textcolor{devilscss}{eher informeller Text mit interessanten Denkanstößen}}
%
\item Douglas Hofstadter:\\ \emph{Gödel, Escher, Bach: An Eternal Golden Braid.} Basic Books\\
{\footnotesize\textcolor{devilscss}{der Klassiker}}
\end{itemize}

\end{frame}


\begin{frame}\frametitle{Logik für InformatikerInnen}

\pause
\begin{center}
\Huge Warum?
\end{center}

\end{frame}

\begin{frame}\frametitle{Logik für InformatikerInnen: Darum (1)}

\begin{center}
\LARGE 
\[ \frac{\text{Logik}}{\text{Informatik}} = \frac{\text{Analysis}}{\text{Ingenieurwesen}}\]

\bigskip
\large
"`Wer rechnende Systeme verstehen und konstruieren will, der benötigt passende mathematische Modelle.\\Dieser Weg führt oftmals zur Logik."'
\end{center}


\begin{itemize}
\item Modellierung von Programm\alert{logik} und \alert{logischen} Schaltungen
\item Berechnung praktisch relevanter Eigenschaften durch logisches Schließen
\item Hauptanwendung: \redalert{Verifikation von Systemen}
\end{itemize}

\end{frame}

\begin{frame}\frametitle{Logik für InformatikerInnen: Darum (2)}

\begin{center}
\LARGE 
Logik = Wissenschaft vom folgerichtigen Denken

\bigskip
\large
"`Wer intelligente Software entwickeln will, der muss logische Schlussfolgerungen algorithmisch umsetzen.\\
Die Logik liefert die nötigen Methoden."'
\end{center}


\begin{itemize}
\item Kodierung von gültigen Zusammenhängen und Regeln
\item Logisches Schließen als Simulation von intelligentem Denken
\item Hauptanwendung: \redalert{Künstliche Intelligenz}
\end{itemize}

\end{frame}


\begin{frame}\frametitle{Logik für InformatikerInnen: Darum (3)}

\begin{center}
\LARGE 
Logisches Schließen = Problemlösen

\bigskip
\large
"`Bedeutende Klassen von (schweren) Problemen lassen sich\\ durch logische Schlussfolgerung lösen.\\
Algorithmen aus der Logik sind in vielen anderen Bereichen anwendbar."'
\end{center}


\begin{itemize}
\item Logik als Spezifikationssprache für komplexe Probleme
\item Logisches Schließen als Suche nach zulässigen Lösungen
\item Anwendungen: \redalert{Constraint-Satisfaction-Probleme} und verwandte Optimierungsaufgaben
\end{itemize}

\end{frame}

\begin{frame}\frametitle{Logik für InformatikerInnen: Darum (4)}

\begin{center}
\LARGE 
Logiken = Beschreibungssprachen

\bigskip
\large
"`Überall wo Informationen maschinell kodiert werden und wo exakt spezifiziert ist, was eine Anwendung\\ aus dieser Kodierung ableiten darf,\\hat man es mit einer Art Logik zu tun."'
\end{center}


\begin{itemize}
\item Logik als Oberbegriff exakt spezifizierter Datenformate
\item Schlussfolgerung zur \ghost{Interpretation/Analyse/Optimierung}
\item Anwendungen: \redalert{Wissensrepräsentation} und \redalert{Datenbanken}
\end{itemize}
\end{frame}

\begin{frame}\frametitle{Was ist Logik?}
\pause

"`Logik"' ist ein allgemeiner Oberbegriff für viele mathematische und technische Formalismen,
gekennzeichnet durch:
\begin{itemize}
\item \redalert{Syntax:} Sprache einer Logik (normalerweise Formeln mit logischen Operatoren)
\item \redalert{Semantik:} Definition der Bedeutung (Worauf beziehen sich die Formeln? Wann ist eine Formel wahr oder falsch?)
% \item \redalert{Inferenz:} Spezifikation der korrekten logische Schlussfolgerungen (Welche neuen Formeln kann man aus den bekannten ableiten?)
\end{itemize}\bigskip\pause

Typische Zielstellung: \redalert{Logische Schlussfolgerung}
\begin{itemize}
\item Welche Schlüsse kann man aus einer gegebenen (Menge von) Formel(n) ziehen?
\item Spezifikation der korrekten Schlussfolgerungen sollte sich aus Semantik ergeben
\item Praktische Berechnung von Schlussfolgerungen ist oft kompliziert
\end{itemize}

\end{frame}

\begin{frame}\frametitle{Viele Logiken}

Es gibt sehr viele Logiken, z.B.
\begin{itemize}
\item Aussagenlogik
\item Prädikatenlogik (erster Stufe)
\item Prädikatenlogik zweiter Stufe
\item Beschreibungslogiken (Wissensrepräsentation und KI, z.B. Modallogik $\mathcal{ALC}$)
\item Temporallogiken (z.B. Modallogik LTL)
\item Logikprogramme (Datalog, Answer Set Programming, Prolog, \ldots)
\item Nichtklassische Logiken (z.B. intuitionistische Logik)
\item Mehrwertige Logiken (z.B. probabilistische Logik)
\item \ldots und viele andere mehr
\end{itemize}

$\leadsto$ In dieser Vorlesung lernen wir zunächst \alert{Aussagenlogik} kennen,\\\phantom{$\leadsto${}} später dann \alert{Prädikatenlogik}, \alert{Modallogiken} und \alert{Datalog}

\end{frame}


\sectionSlide{Aussagenlogik}

\begin{frame}\frametitle{Aussagenlogik}

Die Aussagenlogik untersucht \alert{logische Verknüpfungen von atomaren Aussagen}.
\bigskip

\redalert{Atomare Aussagen} sind Behauptungen, die wahr oder falsch sein können, z.B.:

\examplebox{\vspace{-2ex}
\begin{enumerate}[{A}1]
\item "`Morgen schneit es."'
\item "`Wir werden einen Schneemann bauen."'
% \item "`"'
\end{enumerate}\vspace{-1ex}
\begin{enumerate}[{B}1]
\item "`Die Vorstellung von der globalen Erderwärmung wurde von den Chinesen erfunden."'
\item "`Die Temperatur der Ozeane ist seit 1998 in unverminderter Geschwindigkeit angestiegen."'
\end{enumerate}\vspace{-1ex}
\begin{enumerate}[{C}1]
\item "`Typ-2-Sprachen sind unter Vereinigung abgeschlossen."'
\item "`Typ-2-Sprachen sind unter Komplement abgeschlossen."'
\item "`Typ-2-Sprachen sind unter Schnitt abgeschlossen."'
\end{enumerate}
}

\end{frame}

\begin{frame}\frametitle{Aussagen verknüpfen}

Atomare Aussagen können mithilfe logischer \alert{Junktoren} (oder \alert{Operatoren}) verknüpft werden:

\examplebox{\vspace{-2ex}
\begin{itemize}
\item $\text{A1}\to \text{A2}$: "`Wenn es morgen schneit, dann werden wir einen Schneemann bauen."'
\item $\text{A1}\vee \neg\text{A1}$: "`Entweder schneit es morgen oder nicht."'
\item $\text{B1}\to \neg\text{B2}$: "`Falls die Chinesen die globale Erwärmung erfunden haben, dann steigt die Temperatur der Ozeane nicht unvermindert an."'
\item $(\text{C1}\wedge\text{C2})\to \text{C3}$: "`Sind Typ-2-Sprachen unter Vereinigung und Komplement abgeschlossen, dann sind sie auch unter Schnitt abgeschlossen."'
\end{itemize}}

\end{frame}

% \begin{frame}
% 
% ~\hfill
% \includegraphics[height=6.5cm]{a6}
% \hfill~
% % \rotatebox{90}{\tiny Randall Munroe, \url{http://xkcd.com/208/}, CC-BY-NC 2.5}
% 
% \end{frame}

\begin{frame}\frametitle{Logisches Schließen}

Wenn einige Formeln als wahr angenommen werden, dann kann die Wahrheit anderer Formeln daraus abgeleitet werden.
\bigskip

\examplebox{%
\emph{Beispiele:}
\begin{itemize}
\item Aus $\text{B2}$ und $\text{B1}\to \neg\text{B2}$ folgt $\neg\text{B1}$
\item Aus $\text{C1}$, $\neg\text{C3}$ und $(\text{C1}\wedge\text{C2})\to \text{C3}$ folgt $\neg\text{C2}$
\item Aus $\text{A1}\to \text{A2}$ und $\neg\text{A1}$ folgt \redalert{nicht} $\neg\text{A2}$
\end{itemize}}\pause\bigskip

\anybox{purple}{
Die Gültigkeit bestimmter Schlussfolgerungen hat nichts mit Schneemännern,\\ Erderhitzung oder Typ-2-Sprachen zu tun!\\
Sie ist eine rein logische Konsequenz.
}

\end{frame}

\newcommand{\Aname}{Anna}
\newcommand{\Bname}{Barbara}
\newcommand{\Cname}{Chris}

\begin{frame}\frametitle{Beispiel: Logelei}

\Aname{} behauptet: "`\Bname{} lügt!"'
\bigskip

\Bname{} behauptet: "`\Cname{} lügt!"'
\bigskip

\Cname{} behauptet: "`\Aname{} und \Bname{} lügen!"'
\bigskip\bigskip

{\Large
\narrowcentering{\alert{Wer lügt?}}
}

\medskip
{\footnotesize\narrowcentering{(Und wie kann man das beweisen?)}}
% \bigskip

% 
% $A\leftrightarrow\neg B$
% $B\leftrightarrow\neg C$
% $C\leftrightarrow\neg A \wedge \neg B$

\end{frame}

\begin{frame}\frametitle{Aussagenlogik: Syntax}

Wir betrachten eine abzählbar unendliche Menge $\Slang{P}$ von
\redalert{atomaren Aussagen} (auch bekannt als: \redalert{aussagenlogische Variablen}, \redalert{Propositionen} oder schlicht \redalert{Atome})
\medskip

\defbox{Die Menge der \redalert{aussagenlogischen Formeln} ist induktiv\footnote{Das bedeutet: Die Definition ist selbstbezüglich und soll die kleinste\\ Menge an Formeln beschreiben, die alle Bedingungen erfüllen.} definiert:
\begin{itemize}
\item Jedes Atom $p\in\Slang{P}$ ist eine aussagenlogische Formel
\item Wenn $F$ und $G$ aussagenlogische Formeln sind, so auch:
	\begin{itemize}
	\item $\neg F$: \redalert{Negation}, "`nicht $F$"'
	\item $(F\wedge G)$: \redalert{Konjunktion}, "`$F$ und $G$"'
	\item $(F\vee G)$: \redalert{Disjunktion}, "`$F$ oder $G$"'
	\item $(F\to G)$: \redalert{Implikation}, "`$F$ impliziert $G$"'
	\item $(F\leftrightarrow G)$: \redalert{Äquivalenz}, "`$F$ ist äquivalent zu $G$"'
	\end{itemize}
\end{itemize}
}

Wir verzichten hier oft auf "`aussagenlogisch"' und sprechen z.B. einfach von "`Formeln"'.

\end{frame}

\begin{frame}\frametitle{Beispiele}

Die folgenden Ausdrücke sind aussagenlogische Formeln:
\begin{itemize}
\item $p$
\item $((p\to q)\to p)\to p$
\item $\neg \neg \neg p$
\end{itemize}
\medskip

Die folgenden Ausdrücke sind \redalert{keine} aussagenlogischen Formeln:
\begin{itemize}
\item $p\wedge q \vee r$ (fehlende Klammern)\\
\anybox{strongyellow}{%
\alert{Vereinfachung:} Äußere Klammern dürfen wegfallen, d.h. wir erlauben z.B. $p\to q$ anstatt auf $(p\to q)$ zu bestehen.
}
\item $(p\leftarrow q)$ (Operator $\leftarrow$ undefiniert)
% \item $(p\leftrightarrow q)$ (Operator $\leftrightarrow$ undefiniert)
% \anybox{strongyellow}{%
% \alert{Notation:} Wir erlauben $(p\leftarrow q)$ als alternative Schreibweise für $(q\to p)$ und $(p\leftrightarrow q)$
% als Abkürzung für $((p\to q)\wedge (q\to p))$.}
\item $\overline{p\wedge q}$ (wir verwenden $\overline{\phantom{x}}$ nicht als Negation)
\end{itemize}

\end{frame}

\begin{frame}\frametitle{Teilformeln}

Wir können Formeln als Wörter über (endlichen Teilmengen aus) dem unendlichen Alphabet
$\Slang{P}\cup\{\Sterm{\neg},\Sterm{\wedge},\Sterm{\vee},\Sterm{\to},\Sterm{\leftrightarrow},\Sterm{(},\Sterm{)}\}$ sehen:

\[ \Snterm{F}\to \Slang{P}\mid \Sterm{\neg}\Snterm{F}\mid\Sterm{(}\Snterm{F}\Sterm{\wedge}\Snterm{F}\Sterm{)} \mid\Sterm{(}\Snterm{F}\Sterm{\vee}\Snterm{F}\Sterm{)} \mid\Sterm{(}\Snterm{F}\Sterm{\to}\Snterm{F}\Sterm{)} \mid\Sterm{(}\Snterm{F}\Sterm{\leftrightarrow}\Snterm{F}\Sterm{)}
\]\pause
%
Eine \redalert{Teilformel} ist ein Teilwort (Infix) einer Formel, welches selbst eine Formel ist.
Teilformeln werden auch \redalert{Unterformeln} genannt.\pause
\medskip

Alternativ kann man Teilformeln auch rekursiv definieren:

\defbox{Die \redalert{Menge $\textsf{Sub}(F)$ der Teilformeln} einer Formel $F$ ist definiert als:
\vspace{-2ex}
\[ \textsf{Sub}(F) =\left\{\begin{array}{l@{~~}l@{}}
	\{F\} & \text{falls }F\in\Slang{P} \\
	\{\neg G\}\cup\textsf{Sub}(G) & \text{falls }F=\neg G \\
	\{(G_1\wedge G_2)\}\cup\textsf{Sub}(G_1)\cup\textsf{Sub}(G_2) & \text{falls }F=(G_1\wedge G_2) \\
	\{(G_1\vee G_2)\}\cup\textsf{Sub}(G_1)\cup\textsf{Sub}(G_2) & \text{falls }F=(G_1\vee G_2) \\
	\{(G_1\to G_2)\}\cup\textsf{Sub}(G_1)\cup\textsf{Sub}(G_2) & \text{falls }F=(G_1\to G_2) \\
	\{(G_1\leftrightarrow G_2)\}\cup\textsf{Sub}(G_1)\cup\textsf{Sub}(G_2) & \text{falls }F=(G_1\leftrightarrow G_2)
\end{array}\right.\]\vspace{-2ex}
}

\end{frame}

\begin{frame}\frametitle{Semantik}

\alert{Was bedeutet eine aussagenlogische Formel?}\pause
\begin{itemize}
\item Atome an sich bedeuten zunächst nichts\\
$\leadsto$ sie können einfach wahr oder falsch sein\pause
\item Je nachdem, welche Atome wahr sind und welche falsch, ergeben sich verschiedene "`Interpretationen"'\\
$\leadsto$ dargestellt durch \redalert{Wertzuweisungen}\\
\mbox{}\phantom{$\leadsto$} (Funktionen von $\Slang{P}$ nach $\{\mytrue,\myfalse\}$; $\mytrue$="`wahr"' und $\myfalse$="`falsch"')\pause
\item Die Wahrheit (oder Falschheit) einer Formel ergibt sich aus dem Wahrheitswert der in ihr vorkommenden Atome\\
$\leadsto$ Wertzuweisungen machen Formeln wahr oder falsch
\end{itemize}

\anybox{purple}{
Die Bedeutung einer Formel besteht darin, dass sie uns Informationen darüber liefert,
welche Wertzuweisungen möglich sind, falls die Formel wahr sein soll.
}

\end{frame}

\begin{frame}\frametitle{Wertzuweisungen können Formeln erfüllen}

Eine \redalert{Wertzuweisung} ist eine Funktion $w:\Slang{P}\to\{\mytrue,\myfalse\}$
\bigskip

% Die Wahrheit einer Formel kann rekursiv definiert werden:

\defbox{Eine Wertzuweisung $w$ \redalert{erfüllt} eine Formel $F$, in Symbolen \redalert{$w\models F$}, wenn
eine der folgenden rekursiven Bedingungen gilt:\smallskip

\narrowcentering{\begin{tabular}{rll}
\rowcolor{darkred!70!gray}
\textcolor{white}{Form von $F$} & \textcolor{white}{$w\models F$ wenn:} & \textcolor{white}{$w\not\models F$ wenn:}\\
$F\in \Slang{P}$: & $w(F)=\mytrue$ & $w(F)=\myfalse$\\
\rowcolor{lightred!20}
$F=\neg G$ & $w\not\models G$ & $w\models G$\\
$F=(G_1\wedge G_2)$ & $w\models G_1$ und $w\models G_2$ & $w\not\models G_1$ oder $w\not\models G_2$\\
\rowcolor{lightred!20}
$F=(G_1\vee G_2)$ & $w\models G_1$ oder $w\models G_2$ & $w\not\models G_1$ und $w\not\models G_2$\\
$F=(G_1\to G_2)$ & $w\not\models G_1$ oder $w\models G_2$ & $w\models G_1$ und $w\not\models G_2$\\
\rowcolor{lightred!20}
$F=(G_1\leftrightarrow G_2)$ & $w\models G_1$ und $w\models G_2$ & $w\models G_1$ und $w\not\models G_2$\\%[-1ex]
\rowcolor{lightred!20}
	& \multicolumn{1}{c}{\raisebox{1ex}{oder}} & \multicolumn{1}{c}{\raisebox{1ex}{oder}} \\[-2ex]
\rowcolor{lightred!20}
	& $w\not\models G_1$ und $w\not\models G_2$ & $w\not\models G_1$ und $w\models G_2$\\[-1ex]
\\[-1.5ex]
\end{tabular}}

% Dabei schließt "`oder"' immer den Fall ein, dass beide Möglichkeiten gelten.
Dabei bedeutet "`A oder B"' immer "`A oder B oder beides"'.

% 
% \begin{tabular}{rl}
% $F\in \Slang{P}$: & $w\models F$ gdw. $w(F)=\mytrue$\\
% $F=\neg G$: & $w\models F$ gdw. $w\not\models G$\\
% $F=(G_1\wedge G_2)$: & $w\models F$ gdw. $w\models G_1$ und $w\models G_2$
% \end{tabular}
% 
% \begin{itemize}
% \item $F\in \Slang{P}$ mit $w(F)=\mytrue$
% \item $F=\neg G$ mit $w\not\models G$
% \item $F=(G_1\wedge G_2)$ mit $w\models G_1$ und $w\models G_2$
% \item $F=(G_1\vee G_2)$ mit $w\models G_1$ oder $w\models G_2$ (oder beides!)
% \item $F=(G_1\to G_2)$ mit $w\not\models G_1$ oder $w\models G_2$
% \item $F=(G_1\leftrightarrow G_2)$ mit $w\models G_1$ genau dann wenn $w\models G_2$
% \end{itemize}
% 
% Der \redalert{Wahrheitswert $\hat{w}(F)$} einer Formel $F$ unter einer Wertezuweisung $w$ ist
% wie folgt definiert:
% \vspace{-2ex}
% \[ \hat{w}(F) =\left\{\begin{array}{l@{~~}l@{}}
% 	w(F) & \text{falls }F\in\Slang{P} \\
% 	\{\neg G\}\cup\textsf{Sub}(G) & \text{falls }F=\neg G \\
% 	\{(G_1\wedge G_2)\}\cup\textsf{Sub}(G_1)\cup\textsf{Sub}(G_2) & \text{falls }F=(G_1\wedge G_2) \\
% 	\{(G_1\vee G_2)\}\cup\textsf{Sub}(G_1)\cup\textsf{Sub}(G_2) & \text{falls }F=(G_1\vee G_2) \\
% 	\{(G_1\to G_2)\}\cup\textsf{Sub}(G_1)\cup\textsf{Sub}(G_2) & \text{falls }F=(G_1\to G_2) \\
% 	\{(G_1\leftrightarrow G_2)\}\cup\textsf{Sub}(G_1)\cup\textsf{Sub}(G_2) & \text{falls }F=(G_1\leftrightarrow G_2)
% \end{array}\right.\]\vspace{-2ex}
}

\end{frame}

\begin{frame}\frametitle{Formeln Wahrheitswerte zuweisen}

Wir können Wertzuweisungen von Atomen auf Formeln erweitern:
\[ w(F)=\left\{\begin{array}{rl}\mytrue & \text{falls } w\models F\\
\myfalse & \text{falls } w\not\models F
\end{array}\right.\]\pause

\redalert{Wahrheitswertetabellen} illustrieren die Semantik der Junktoren:\medskip

\narrowcentering{%
\begin{tabular}{cc}
\rowcolor{lightblue!20}
$w(F)$ & $w(\neg F)$\\
\myfalse & \mytrue \\
\rowcolor{gray!10}
\mytrue & \myfalse \\
\end{tabular}}\medskip

\narrowcentering{\begin{tabular}{cccccc}
\rowcolor{lightblue!20}
$w(F)$ & $w(G)$ & $w(F\wedge G)$ & $w(F\vee G)$ & $w(F\to G)$ & $w(F\leftrightarrow G)$\\
\myfalse & \myfalse & \myfalse & \myfalse& \mytrue & \mytrue\\
\rowcolor{gray!10}
\mytrue & \myfalse & \myfalse & \mytrue& \myfalse & \myfalse\\
\myfalse & \mytrue & \myfalse & \mytrue& \mytrue & \myfalse\\
\rowcolor{gray!10}
\mytrue & \mytrue & \mytrue & \mytrue& \mytrue & \mytrue\\
\end{tabular}}

\end{frame}

\begin{frame}\frametitle{Wahrheitswerte von Formeln bestimmen}

\begin{itemize}
\item Der Wahrheitswert einer Formel hängt nur vom Wahrheitswert der (endlich vielen) Atome ab, die in
ihr vorkommen.\\
$\leadsto$ Wir geben oft nur diese an.
\item Der Wahrheitswert einer Formel ergibt sich rekursiv aus dem Wahrheitswert ihrer Teilformeln.\\
$\leadsto$ Darstellung in Wahrheitswertetabelle
\end{itemize}

\examplebox{\emph{Beispiel:} Für die Formel $F=((p\to q)\leftrightarrow(\neg q\to\neg p))$ und Wertzuweisung 
$w$ mit $w(p)=\myfalse$ und $w(q)=\mytrue$ ergibt sich die folgende Tabelle:
\medskip

\narrowcentering{\begin{tabular}{ccccccc}
\rowcolor{darkgreen!30}
$w(p)$ & $w(q)$ & $w(p\to q)$ & $w(\neg p)$ & $w(\neg q)$ & $w(\neg q\to \neg p)$ & $w(F)$\\
\myfalse & \mytrue & \mytrue & \mytrue & \myfalse & \mytrue & \mytrue\\
\end{tabular}}

}

\end{frame}

\begin{frame}\frametitle{Beispiel: Logelei}

\Aname{} behauptet: "`\Bname{} lügt!"'\hfill \visible<2->{$A\leftrightarrow \neg B$}
\bigskip

\Bname{} behauptet: "`\Cname{} lügt!"'\hfill \visible<3->{$B\leftrightarrow \neg C$}
\bigskip

\Cname{} behauptet: "`\Aname{} und \Bname{} lügen!"'\hfill \visible<4->{$C\leftrightarrow (\neg A\wedge \neg B)$}
\bigskip\bigskip

{\Large
\narrowcentering{\alert{Wer lügt?}}
}

\medskip
{\footnotesize\narrowcentering{(Und wie kann man das beweisen?)}}
% \bigskip

% 
% $A\leftrightarrow\neg B$
% $B\leftrightarrow\neg C$
% $C\leftrightarrow\neg A \wedge \neg B$

\end{frame}

% \begin{frame}\frametitle{Logische Konsequenzen}
% 
% \defbox{Eine Wertzuweisung $w$ ist \redalert{Modell einer Formel $F$} wenn $w\models F$ (also wenn $w(F)=\mytrue$).
% % Wir erweitern diese Notation auf (möglicherweise unendliche) Formelmengen: $w$ ist ein \redalert{Modell der Menge $\mathcal{F}$}
% %  wenn $w\models F$ für alle $F\in\mathcal{F}$ gilt. In diesem Fall schreiben wir \redalert{$w\models\mathcal{F}$}.
% % 
% Ist $\mathcal{F}$ eine (möglicherweise unendliche) Menge von Formeln, dann ist $w$ ein \redalert{Modell der Menge $\mathcal{F}$}
%  wenn $w\models F$ für alle $F\in\mathcal{F}$ gilt. In diesem Fall schreiben wir \redalert{$w\models\mathcal{F}$}.
% }\pause
% 
% Die logischen Schlussfolgerungen aus einer Formel(menge) ergeben sich aus ihren Modellen:
% 
% \defbox{Sei $\mathcal{F}$ eine Menge von Formeln.
% Eine Formel $G$ ist eine \redalert{logische Konsequenz} aus $\mathcal{F}$ wenn jedes Modell
% von $\mathcal{F}$ auch ein Modell von $G$ ist.\\ In diesem Fall schreiben wir $\mathcal{F}\models G$.
% }
% 
% \begin{enumerate}[$\leadsto$]
% \item Die Formeln in $\mathcal{F}$ schränken die möglichen Interpretationen \ghost{ein:}\\
% Je mehr Formeln wahr sein sollen,\\ desto weniger Freiheiten gibt es bei der Wahl der Modelle
% \item $G$ ist eine logische Konsequenz wenn gilt:\\ Falls $\mathcal{F}$ wahr ist, dann ist auch $G$ garantiert wahr
% % 
% % Eine logische Konsequenz $G$ von $\mathcal{F}$ ist in jedem Fall (jedem Modell wahr in dem $\mathcal{F}$ wahr ist
% \end{enumerate}
% 
% \end{frame}
% 
% \begin{frame}[t]\frametitle{Beispiel: Logelei}
% 
% \Aname{} behauptet: "`\Bname{} lügt!"'\hfill {$A\leftrightarrow \neg B$}\\[1ex]
% \Bname{} behauptet: "`\Cname{} lügt!"'\hfill {$B\leftrightarrow \neg C$}\\[1ex]
% \Cname{} behauptet: "`\Aname{} und \Bname{} lügen!"'\hfill {$C\leftrightarrow (\neg A\wedge \neg B)$}\\[2ex]
% 
% \begin{tabular}{c@{~ }c@{~ }cc@{\hspace{3mm}}c@{\hspace{3mm}}c}
% \rowcolor{lightblue!20}
% $w(A)$ & $w(B)$ & $w(C)$ & $w(A\leftrightarrow \neg B)$ & $w(B\leftrightarrow \neg C)$ & $w(C\leftrightarrow (\neg A\wedge \neg B))$\\
% \myfalse & \myfalse & \myfalse & \myfalse& \myfalse & \myfalse\\
% \rowcolor{gray!10}
% \mytrue & \myfalse & \myfalse & \mytrue& \myfalse & \mytrue\\
% \myfalse & \mytrue & \myfalse & \mytrue& \mytrue & \mytrue\\
% \rowcolor{gray!10}
% \mytrue & \mytrue & \myfalse & \myfalse& \mytrue & \mytrue\\
% \myfalse & \myfalse & \mytrue & \myfalse& \mytrue & \mytrue\\
% \rowcolor{gray!10}
% \mytrue & \myfalse & \mytrue & \mytrue& \mytrue & \myfalse\\
% \myfalse & \mytrue & \mytrue & \mytrue& \myfalse & \myfalse\\
% \rowcolor{gray!10}
% \mytrue & \mytrue & \mytrue & \myfalse& \myfalse & \myfalse\\
% \end{tabular}
% 
% % {\Large
% % \narrowcentering{\alert{Wer lügt?}}
% % }
% % 
% % \medskip
% % {\footnotesize\narrowcentering{(Und wie kann man das beweisen?)}}
% % \bigskip
% 
% 
% \end{frame}
% 
% \begin{frame}[t]\frametitle{Beispiel: Logelei (2)}
% 
% \Aname{} behauptet: "`\Bname{} lügt!"'\hfill {$A\leftrightarrow \neg B$}\\[1ex]
% \Bname{} behauptet: "`\Cname{} lügt!"'\hfill {$B\leftrightarrow \neg C$}\\[1ex]
% \Cname{} behauptet: "`\Aname{} und \Bname{} lügen!"'\hfill {$C\leftrightarrow (\neg A\wedge \neg B)$}\\[2ex]
% 
% \begin{tabular}{c@{~ }c@{~ }cc@{\hspace{3mm}}c@{\hspace{3mm}}c}
% \rowcolor{lightblue!20}
% $w(A)$ & $w(B)$ & $w(C)$ & $w(A\leftrightarrow \neg B)$ & $w(B\leftrightarrow \neg C)$ & $w(C\leftrightarrow (\neg A\wedge \neg B))$\\
% % \myfalse & \myfalse & \myfalse & \myfalse& \myfalse & \myfalse\\
% % \rowcolor{gray!10}
% % \mytrue & \myfalse & \myfalse & \mytrue& \myfalse & \mytrue\\
% \myfalse & \mytrue & \myfalse & \mytrue& \mytrue & \mytrue\\
% % \rowcolor{gray!10}
% % \mytrue & \mytrue & \myfalse & \myfalse& \mytrue & \mytrue\\
% % \myfalse & \myfalse & \mytrue & \myfalse& \mytrue & \mytrue\\
% % \rowcolor{gray!10}
% % \mytrue & \myfalse & \mytrue & \mytrue& \mytrue & \myfalse\\
% % \myfalse & \mytrue & \mytrue & \mytrue& \myfalse & \myfalse\\
% % \rowcolor{gray!10}
% % \mytrue & \mytrue & \mytrue & \myfalse& \myfalse & \myfalse\\
% \end{tabular}
% 
% $\leadsto$ Genau ein Modell (bezüglich der relevanten Atome)
% \pause\bigskip
% 
% \emph{Logische Konsequenzen:} Alle Formeln, die unter einer Wertzuweisung mit $w(A)=\myfalse$, $w(B)=\mytrue$ und $w(C)=\myfalse$ wahr sind, \ghost{z.B.:}
% \begin{itemize}
% \item $\neg A$ ("`\Aname{} lügt"')
% \item $\neg C$ ("`\Cname{} lügt"')
% \item $\neg A\wedge \neg C$ ("`\Aname{} und \Cname{} lügen"')
% \item $B$ ("`\Bname{} sagt die Wahrheit"')
% \item Aber auch: $D\vee \neg D$
% \end{itemize}
% 
% % {\Large
% % \narrowcentering{\alert{Wer lügt?}}
% % }
% % 
% % \medskip
% % {\footnotesize\narrowcentering{(Und wie kann man das beweisen?)}}
% % \bigskip
% 
% 
% \end{frame}
% 
% 
% \begin{frame}\frametitle{Allgemeingültigkeit und Co.}
% 
% \defbox{Eine Formel $F$ ist:
% \begin{itemize}
% \item \redalert{unerfüllbar} (oder \alert{inkonsistent}), wenn sie keine Modelle hat
% \item \redalert{erfüllbar} (oder \alert{konsistent}), wenn sie Modelle hat
% \item \redalert{allgemeingültig} (oder eine \alert{Tautologie}), wenn alle Wertzuweisungen Modelle für die Formel sind
% \item \redalert{widerlegbar}, wenn sie nicht allgemeingültig ist
% \end{itemize}
% Diese Begriffe kann man für Mengen von Formeln genauso definieren.
% }
% 
% \examplebox{\emph{Beispiel:} 
% \begin{itemize}
% \item $p\wedge\neg p$ ist unerfüllbar (und widerlegbar)
% \item $p\vee\neg p$ ist allgemeingültig (und erfüllbar)
% \item $p\wedge q$ ist erfüllbar und widerlegbar
% \end{itemize}}
% 
% \end{frame}
% 
% \begin{frame}\frametitle{Allgemeingültigkeit und Unerfüllbarkeit}
% 
% % Speziell Allgemeingültigkeit und Unerfüllbarkeit sind für logische Schlussfolgerung von Bedeutung:
% 
% \theobox{\emph{Satz:}
% \begin{enumerate}[(1)]
% \item Eine allgemeingültige Formel ist logische Konsequenz jeder anderen Formel(menge).
% \item Eine unerfüllbare Formel(menge) hat jede andere Formel als logische Konsequenz.
% \end{enumerate}}
% 
% \emph{Beweis:} Folgt direkt aus Definition von logischer Konsequenz. Für (2) ist es wichtig, dass eine Eigenschaft "`für alle Modelle"' gilt, wenn es keine Modelle gibt (sie gilt dann für "`alle null Modelle"'). \qed
% \bigskip\pause
% 
% \anybox{purple}{
% Der unerwartete(?) Effekt (2) ist im Deutschen sprichwörtlich:\medskip
% 
% \narrowcentering{"`Wenn das stimmt, dann bin ich der Kaiser von China!"'}\medskip
% 
% drückt aus, dass aus einer mutmaßlich falschen Annahme alles folgt, selbst wenn 
% es offensichtlich unwahr ist.
% }
% 
% \end{frame}



\begin{frame}\frametitle{Zusammenfassung und Ausblick}

\redalert{Logik}, logisches Modellieren und logisches Schlussfolgern hat in der Informatik viele direkte und indirekte Anwendungen
\bigskip

Mithilfe der \redalert{Aussagenlogik} kann man logische Beziehungen atomarer Aussagen spezifizieren
\bigskip

\redalert{Wertzuweisungen} können eine aussagenlogische Formel erfüllen -- dann nennt man sie \redalert{Modell} -- oder widerlegen
\bigskip

% Die Modelle einer Formel(menge) definieren ihre \redalert{logischen Konsequenzen}
% (="`alles, was in diesen Fällen noch gilt"')\bigskip

\anybox{yellow}{
Offene Fragen:
\begin{itemize}
\item Geht logisches Schließen auch ohne Wahrheitswertetabellen?
\item Wie (in)effizient ist logisches Schließen?
\item Was hat das mit Komplexität zu tun?
\end{itemize}
}

\end{frame}



\end{document}
